\documentclass[a4paper]{article}
    \usepackage{fullpage}
    \usepackage{amsmath}
    \usepackage{amssymb}
    \usepackage{textcomp}
    \usepackage[utf8]{inputenc}
    \usepackage{hyperref}
    \usepackage[T1]{fontenc}
    \textheight=10in
    \pagestyle{empty}
    \raggedright
    \usepackage[left=0.8in,right=0.8in,bottom=0.8in,top=0.8in]{geometry}

\def\bull{\vrule height 0.8ex width .7ex depth -.1ex }

% DEFINITIONS FOR RESUME %%%%%%%%%%%%%%%%%%%%%%%

\newcommand{\lineunder} {
    \vspace*{-8pt} \\
    \hspace*{-18pt} \hrulefill \\
}

\newcommand{\header} [1] {
    {\hspace*{-18pt}\vspace*{6pt} \textsc{#1}}
    \vspace*{-6pt} \lineunder
}

\newcommand{\contact} [3] {
    \vspace*{20pt}
    \begin{center}
        {\Huge \textbf{#1}}\\
        \vspace{5pt}
        #2 \\ #3
    \end{center}
    \vspace*{-8pt}
}

\usepackage{xcolor}
\definecolor{cvblue}{RGB}{0, 50, 140} % Professional Navy

\usepackage{hyperref}
\hypersetup{
    colorlinks=true,
    linkcolor=black,     % Keep internal page links black
    urlcolor=cvblue,      % Make external evidence links blue
    pdfborder={0 0 0}
}

% END RESUME DEFINITIONS %%%%%%%%%%%%%%%%%%%%%%%

\begin{document}
\vspace*{-40pt}

%==== Profile ====%
\vspace*{-10pt}
\begin{center}
    {\Huge \textbf{Aaromal A.}}\\
    \vspace{5pt}
    Kochi, India $\cdot$ aaromalonline@gmail.com $\cdot$ \href{https://www.linkedin.com/in/aaromalonline/}{https://www.linkedin.com/in/aaromalonline/} \\
    \vspace{4pt}
    Self-taught Programmer $\cdot$ FOSS (GNU/Linux) $\cdot$ Electronics Engineering
\end{center}
\vspace*{-8pt}

%==== Education ====%
\vspace*{2mm}
\header{Education}
\textbf{\href{https://www.cbse.gov.in/}{Central Board of Secondary Education (CBSE)}} \hfill Kerala, 2022 -- 2024\\
\textit{Higher Secondary Education, Mathematics and Computer Science} \hfill Grade: 94.6\%\\
\textbf{Coursework:} Python Programming, Database Management Systems (DBMS), Data Structures and Algorithms (DSA), Computer Networking, Statistics & Probability, Coordinate Geometry, Calculus.\\
\textbf{Activities:} Kerala State Junior Meet Athlete, District Junior Athletic Meet Medalist, School Cadre Group Leader.
\vspace{4mm}

\textbf{\href{https://www.cusat.ac.in/}{Cochin University of Science \& Technology (CUSAT)}} \hfill Kerala, 2024 -- 2028\\
\textit{B.Tech in Electronics \& Communication Engineering} \hfill CGPA: 8.96 (S2)\\
\textbf{Major Coursework:} Analog Integrated Circuits, Digital System Design, Signals and Systems, Network Theory, Microprocessors \& Microcontrollers (8086, 8051), C programming \& Object-Oriented Programming (C++), Linear Algebra \& Transform Techniques, Numerical \& Statistical Methods, Calculus.
\vspace*{4mm}

%==== Research & Presentations ====%
\header{Research \& Presentations}
\textbf{\href{https://github.com/aaromalonline/ecometachip}{Ecometachip Mini: Low-Cost Dielectric Chemical Identifier}} \hfill \textit{New Delhi} \\
\textit{\href{https://epapers2.org/apscon2026/ESR/paper_details.php?paper_id=6101}{Lecture Presentation, IEEE APSCON 2026 (Student Research Forum)}} \hfill \textit{Jan 2026}
\begin{itemize} \itemsep 1pt
    \item Designed a dielectric sensing system for preliminary identification of common liquids (water, ethanol, acetone, petrol) based on relative permittivity variation.
    \item Integrated a biodegradable coir-rubber composite as the primary sensing element.
    \item Developed a 100 MHz Hartley oscillator-based readout circuit featuring RF pickup and rectification.
    \item Implemented LED-based visual classification to validate output voltage variations.
    \item Simulated ideal behavior in \textbf{LTspice} by varying coupling constants and analyzed the sensitivity and repeatability of the sensor under variable environmental conditions (Humidity/Temp) using \textbf{Python (Pandas/Seaborn)}.
    \item Quantified environmental drift coefficients ($mV/\%RH$ and $mV/^\circ C$) to characterize sensor robustness and output consistency across variable ambient conditions.
    \item Prepared a detailed research paper using \LaTeX\ which got selected for IEEE APSCON'26.
\end{itemize}
\vspace*{4mm}

%==== Featured Projects ====%
{\hspace*{-18pt}\vspace*{6pt} \textsc{Featured Projects} \hfill \href{https://github.com/aaromalonline}{github.com/aaromalonline}}
\vspace*{-6pt} \lineunder

\textbf{\href{https://liberate-atcs.vercel.app/}{Liberate:ATCS - Adaptive Typing \& Control System}} \hfill \textit{Feb 2025 - Sep 2025}\\
{\sl ESP32, C++, Python, Arduino IDE}
\begin{itemize} \itemsep 1pt
    \item Developed a muscle-twitch-based communication system translating subtle facial/body movements into digital control signals.
    \item Enabled hands-free operation of computer cursors, wheelchairs, and keyboards for individuals with limited mobility.
    \item Received participation recognition at the \textbf{2025 SSCS Arduino Content}.
\end{itemize}
\vspace*{1mm}

\href{https://github.com/aaromalonline/cardiokey}{\textbf{CardioKey (ECGID) - Analog ECG Acquisition \& Biometrics}} \hfill \textit{Feb 2026 - Present}\\
{\sl AD620, ESP32, Python/MATLAB, Scikit-learn}
\begin{itemize} \itemsep 1pt
    \item Engineered an analog front-end using the AD620 instrumentation amplifier to acquire low-amplitude (~1mV) heart signals.
    \item Implemented a P-QRS-T complex extraction pipeline using a 16-bit ADC (ADS1115) and ESP32 for digital signal processing.
    \item Utilized Principal Component Analysis (PCA) to compress 250-point cardiac windows into 30-point biometric "signatures".
    \item Achieved a 96\% identification accuracy using Linear Discriminant Analysis (LDA) and majority vote classifiers.
    \item Featured inherent "Liveness Detection" to prevent biometric spoofing common in fingerprint or face ID systems.
\end{itemize}
\vspace*{2mm}

\href{https://github.com/aaromalonline/PassCrypt}{\textbf{PassCrypt - Password Manager}} \hfill \textit{2022}\\
{\sl Python, SQLite3}
\begin{itemize}
    \item Developed a robust password management application utilizing \textbf{SQLite} for secure local database storage.
    \item Implemented secure \textbf{User Authentication} and access control using advanced hashing techniques (\textbf{SHA256}) to protect sensitive user tables.
    \item Integrated \textbf{Cryptographic protocols} like \textbf{AES} using cryptography.fernet to ensure that all stored passwords remain encrypted and secure from unauthorized access.
    \item Designed an \textbf{Advanced Filtering and Search} interface, enabling users to efficiently sort and manage large lists of credentials.
    \item Built a \textbf{Data Portability} feature allowing users to navigate entries and export/import entire password databases to a single file.
\end{itemize}
\vspace*{2mm}

\href{https://github.com/aaromalonline/truevote}{\textbf{Truevote: Digital Voting System}} \hfill {\sl C++, Qt, SHA256 Hashing, CMake} \\
Secure e-voting replacement for EVMs using OOP principles and SHA256 hashing to ensure data integrity.

\vspace*{4mm}

%==== Experience ====%
\header{Experience}
\textbf{Software Intern} \hfill \textit{Mar 2025 -- Jul 2025} \\
\textit{\href{https://roverchallenge.eu/team/horizon-2024/}{Horizon Mars Rover Team, CUSAT}} \hfill \textit{Kochi, India}
\begin{itemize}
    \item Architected node communication for a functional Mars rover using \textbf{ROS 2 Humble} and \textbf{Arduino programming}.
    \item Performed circuit simulation using \textbf{TinkerCAD} and \textbf{Arduino} to program robotic arm joints using servo-motors.
    \item Wrote technical reports of the rover for upcomming IROC and ERC competitions using \LaTeX\
    \item Track record for qualifying AIR 26 on IRC26' and global rank 18 on ERC24' rover challenges including the \textbf{European Rover Challenge (ERC)} and \textbf{IROC by ISRO}, India.
\end{itemize}

\textbf{Active Member} \hfill \textit{Mar 2025 -- Present} \\
\textit{\href{https://fossunited.org/c/school-of-engineering-cusat}{FOSS CUSAT | Community}} \hfill \textit{Kochi, India}
\begin{itemize}
    \item Spearheaded the \textbf{Linux Installation Party} as part of the \textbf{Endof10 Campaign}, facilitating the migration and onboarding of 20+ students to open-source operating systems.
    \item Provided hands-on technical guidance on \textbf{manual disk partitioning}, Secure Boot configurations, and shell workflows across distributions including Debian, Arch, Fedora, and Pop!\_OS.
    \item Demonstrated proficiency in deployment tools like \textbf{Ventoy} and \textbf{Balena Etcher} to provision Linux environments on various hardware including external SSDs and live USBs.
    \item Represented FOSS CUSAT at major networking events, including \textbf{Kochi FOSS meetups} and \textbf{IndiaFOSS 2025} (Bengaluru), contributing to the regional growth of the free software community.
\end{itemize}
\vspace*{4mm}

\header{Technical Skills \& Certifications}
\begin{itemize} \itemsep 1pt
    \item \textbf{Languages:} Proficient in \textbf{Python} (NumPy, SciPy, Matplotlib, Seaborn, Scikit-learn, Pandas, OpenCV, Requests, Selenium, SQLite), \textbf{C/C++}, HTML/CSS, JavaScript, Bash, SQL, \LaTeX\, and Verilog.
    \item \textbf{Tools \& Frameworks:} \textbf{ROS 2 Humble} (Basic Nodes), \textbf{Qt}, \textbf{Arduino} Framework, \textbf{LTspice}, \textbf{Xilinx Vivado}, Proteus, and \textbf{MATLAB \& Simulink} (Basics).
    \item \textbf{Linux \& DevOps:} \textbf{Git} workflow; Daily-driving \textbf{Linux} for $>$5 years across \textbf{Debian} (current), Ubuntu, Fedora, and Arch; experienced in VM management, manual SSD/HDD partitioning, Ventoy, and Balena Etcher.
    \item \textbf{MATLAB Academy (Online Summer Internship):} \hfill \textit{Jun 2025}
    \begin{itemize}
        \item \href{https://www.credly.com/badges/5c9e05db-f4b5-4435-a360-226a04411a5e}{\textbf{MATLAB Proficiency:}} Core programming, algorithm development, and signal processing.
        \item \href{https://www.credly.com/badges/47d8690a-7a2d-4248-9c41-27fc8f8edc62}{\textbf{Data Analysis \& Processing:}} Visualization and AI/Machine Learning modules.
        \item \href{https://www.credly.com/badges/f69c1b8e-5a25-423f-9eba-db02fd2b6fb2}{\textbf{MATLAB for Simulink Users:}} System modeling, circuit simulation, and control logic.
    \end{itemize}
\end{itemize}

\end{document}
